\section{Cosa è il Modello delle classi di progetto?}
Il diagramma delle classi di progetto (DCD - Design Class Diagram)  è un modello utilizzato per la modellazione statica degli oggetti, e sserve a illustrare le classi, le interfacce e le relative associazioni adottando uno specifico punto di vista software
\subsection{Differenza col Modello di Dominio}
A differenza del \textbf{Modello di Dominio}, \uline{che rappresenta gli elementi da un punto di vista concettuale e slegato dal codice}, il \textbf{Modello delle Classi di Progetto} \uline{definisce come il sistema sarà realmente implementato}
\section{Parole Chiave}
\begin{tabularx}{\linewidth}{|l|X|X|}
\hline \textbf{Parola Chiave} & \textbf{Significato} & \textbf{Esempio di suo} \\ \hline
\texttt{<<actor>>} & il classificatore è un attore & Nei diagrammi delle classi, sopra al nome di un classificatore \\ \hline
\texttt{<<interface>>} & Il classificatore è un'interfaccia & Nei diagrammi delle classi, sopra al nome di un classificatore \\ \hline
\texttt{\{abstract\}} & l'elemento è astratto, non può essere instanziato & Nei diagrammi delle classi, dopo il nome di un classificatore o il nome di un'operazione \\ \hline
\texttt{\{ordered\}} & Un insieme di oggetti ha un ordine predefinito & Nei diagrammi delle classi, a un'estremità di associazione \\ \hline
\end{tabularx}
\section{Componenti del DCD}
\subsection{Classificato (Classi e interfacce)}
\subsubsection{Classi}
Rappresentano la struttura portante e modellano un insieme di oggetti con attributi e operazioni comuni
\begin{figure}[H]
    \centering
    \includegraphics[scale=0.6]{Schermata_20260429_113757.png}
    \label{Rappresentazione grafica di una classe}
\end{figure}
\subsection{Interfacce}
Definiscono un "contratto" (un insieme di operazioni pubbliche senza implementazione) che le classi devono rispettare.
\subsection{Attributi}
Gli attributi modella le proprietà locali e interne di una classe

In un DCD, a differenza del modello di dominio concettuale, presentano dettagli tecnici specifici
\subsubsection{Sintassi}
Includono il nome, il tipo di dato e l'eventuale molteplicità o valore iniziale
\subsubsection{Visibilità}
\begin{tabularx}{\linewidth}{|X|X|X|}
\hline
\textbf{Simbolo} & \textbf{Nome} & \textbf{Semantica} \\ \hline
+ & Pubblico & Visibile a tutti \\ \hline
- & Privato & Visibile solo alla classe \\ \hline
\# & Protetta & Visibile alla classe e ai suoi "discendenti" \\ \hline
\textasciitilde & Package & Visibile ad ogni elemento dello stesso package \\ \hline
\end{tabularx}
\subsection{Operazioni e Metodi}
\subsubsection{Operazioni}
Sono le dichiarazioni della firma
\subsubsection{Metodi}
Sono l'implementazione pratica dell'operazione
\subsection{Associazioni}
Le associazioni vengono disegnate come linee che uniscono le classi e si definiscono attraverso diversse proprietà 
\subsubsection{Nome e verso}
Si può assegnare un nome descrittivo all'asociazione.

Talvolta si accompagna a un fraccia (il verso) che aiuta solo a leggere l'associazione nel senso corretto

\begin{figure}[H]
    \centering
    \includegraphics[scale=0.65]{Schermata_20260429_121526.png}
\end{figure}

\subsubsection{Ruoli}
Si indicano scrivendo un nome all'estremità della linea per specificare quale ruola gioca quella classe nell'associazione.
\begin{figure}[H]
    \centering
    \includegraphics[scale=0.65]{Schermata_20260429_121155.png}
\end{figure}
\subsubsection{Molteplicità}
Sono dei limiti numerici posti vicino alla destinazione che indicano quanti oggetti di quela classe possono partecipare alla relazione in un dato istante
\begin{figure}[H]
    \begin{figure}[H]
    \centering
    \includegraphics[scale=0.65]{Schermata_20260429_121919.png}
\end{figure}
\end{figure}
\subsection{Associazioni n-arie}
